% Generated draft: requires mathematical and source review.
\documentclass[11pt]{article}
\usepackage[margin=25mm]{geometry}
\usepackage{fontspec}
\setmainfont{Latin Modern Roman}
\usepackage{amsmath,amssymb,mathtools,unicode-math}
\setmathfont{Latin Modern Math}
\usepackage{enumitem,needspace,xurl,hyperref}
\setlength{\parindent}{0pt}
\setlength{\parskip}{5pt}
\setlength{\emergencystretch}{3em}
\newcommand{\N}{\mathbb{N}}
\newcommand{\Z}{\mathbb{Z}}
\newcommand{\Q}{\mathbb{Q}}
\newcommand{\R}{\mathbb{R}}
\newcommand{\C}{\mathbb{C}}
\newcommand{\F}{\mathbb{F}}
\newcommand{\A}{\mathbb{A}}
\newcommand{\PP}{\mathbb{P}}
\newcommand{\E}{\mathbb{E}}
\newcommand{\eps}{\varepsilon}
\newcommand{\dd}{\,\mathrm{d}}
\begin{document}
\section*{UnsolvedMath ID: 30002233}
\textbf{Source problem code:} catalogue entry
\subsection{Definitions and mathematical statement}
Let \(K\) be the field appearing in the supplied record, and let
\[
f\colon X'\longrightarrow X
\]
be a finite morphism of smooth, proper, geometrically connected curves over \(K\). The record specifies an induced map
\[
f_*\colon\Phi_{X'}\longrightarrow\Phi_X,
\]
but does not define the groups \(\Phi_{X'}\) and \(\Phi_X\) or the construction of this map.

The supplied question asks whether, for every such morphism satisfying the additional conditions below, \(f_*\) is surjective:
\begin{enumerate}
\item the minimal regular model of \(X'\) ``consists of two rational components meeting transversally'';
\item \(X\) is a Mumford curve and \(g(X)\geq 2\), where \(g(X)\) denotes its genus.
\end{enumerate}
Here surjectivity means
\[
\forall\phi\in\Phi_X\quad\exists\phi'\in\Phi_{X'}\quad f_*(\phi')=\phi.
\]

\textit{Unresolved qualifications.} This is a schematic rendering of the supplied question, not a fully specified conjecture. The record gives no valuation, completeness, residue-field, or other hypotheses on \(K\), and supplies no base for the regular model. It is unclear whether the condition on two components refers to the special fiber of that model; this interpretation has not been silently substituted for the supplied wording. The field over which the components are rational, their multiplicities, and their number of intersection points are unspecified. The meaning of Mumford curve over the intended field, including any splitting conventions, also requires clarification. A possible interpretation of \(\Phi_Y\) is the component group of the special fiber of the N\'eron model of the Jacobian of \(Y\), with \(f_*\) induced by pushforward on Jacobians, but the supplied record does not establish that interpretation or whether surjectivity concerns group schemes, geometric component groups, or residue-field rational points.

The supplied dated assessment additionally says that the extracted record is only the lead-in to Question 21 and contains no independently stated mathematical claim. Its relationship to the explicit question reproduced here must therefore be checked before this draft is treated as an authoritative formulation.
\subsection{Short English statement}
The supplied question asks whether a finite morphism of curves induces a surjection on the indicated component groups when the source has the stated two-component model and the target is a Mumford curve of genus at least \(2\). The field, component groups, model condition, and source context need clarification. The supplied open status is a curation label, not an independently verified mathematical conclusion.
\subsection{Sources}
\begin{itemize}
\item Local UnsolvedMath registry: permanent ID \(30002233\), source problem code \texttt{catalogue entry}, title \emph{Surjectivity on Component Groups of Semistable Curves}, in the Oberwolfach Reports Open Problems collection. No individual proposer is supplied; the proposed year is recorded as 2012.
\item Supplied source citation: \emph{Non-Archimedean Analytic Geometry} (2013), Oberwolfach Reports, pp.\ 3211--3265, DOI \texttt{10.4171/owr/2012/53}. Source URL: \url{https://doi.org/10.4171/owr/2012/53}. The background describes this report as documentation of a mathematical workshop and its open-problem discussions.
\item Supplied literature source, labelled ``Oberwolfach report, source context'': \url{https://doi.org/10.4171/owr/2012/53}.
\item Literature assessment dated \texttt{2026-08-21}: the supplied curation labels the record open, but says that the extracted text is only the lead-in to Question 21 and contains no independently stated mathematical claim; it retains the open label under its conservative rule. This is supplied dated curation, not independently verified literature assessment. The background explicitly requests independent verification of time-sensitive status and cited claims.
\item The supplied statement status is \texttt{corrected\_verified}; its accompanying verification note reports manual semantic, typographic, and LaTeX review and claims that the cleanup is self-contained and introduces no unsupported source-dependent content. This is a supplied review assertion, not independent source verification by this draft.
\end{itemize}
\end{document}
